\section{HyMEM：电脑操作智能体的混合自演化记忆}
\label{sec:hymem}

HyMEM\upcite{zhu-etal-2026-hybrid}面向需要观察截图并连续操作网页的电脑操作智能体。历史经验既包含可复用策略，也包含界面与动作细节；任务推进后，最初检索的经验可能失效。系统因此结合离散策略与连续轨迹表示，分别更新长期记忆和当前任务的工作记忆。

\subsection{问题背景与混合表示}

输入包括任务描述、当前截图及历史成功操作轨迹，输出为下一步可执行操作。例如，搜索商品和结账所需经验不同，但可能需要保留同一预算约束。纯文本总结容易丢失视觉细节，单纯连续表示则缺少清晰的策略组织和可解释索引；一次性检索也难以覆盖交互阶段变化。这些问题构成混合表示与两种更新机制的动机。

论文在图 $G=(V,E)$ 中组织经验，将一条轨迹的逻辑表示写为
\begin{equation}
v_i=(c_i,A_i,m_i),
\end{equation}
其中 $c_i$ 为高层策略，$A_i$ 为动作、界面元素或领域标签等属性集合，$m_i$ 为连续轨迹表示。共享相同属性的轨迹间建立无向边，并通过策略、属性与轨迹的关联支持跨经验检索。策略与属性提供语义锚点；借鉴CoMEM的连续编码保留难以用文本概括的视觉及动作细节，每条轨迹压缩为8个嵌入向量并接入智能体输入。

检索使用独立的索引向量：归一化的任务文本CLIP表示与起始截图CLIP表示拼接后，通过FAISS寻找相似经验。连续记忆$m_i$则负责向智能体提供检索到的轨迹细节。

\subsection{全局记忆演化与结构化检索}

更新长期图时，新轨迹先检索相关经验，再由视觉语言模型判断信息增量。检索提供候选集合，模型进一步判断经验是否重复。三种更新为：
\begin{itemize}
    \item \textbf{ADD}：出现新策略或属性，增加对应节点与连接。
    \item \textbf{MERGE}：已有策略获得互补证据，如界面变体或另一条操作路径，将证据并入相应策略及属性关联。
    \item \textbf{REPLACE}：在同等条件下，新轨迹以更少步骤、更高成功率或更低风险实现相同策略，用其替换过时证据并刷新关联。
\end{itemize}
这些操作在长期图中积累可迁移经验并减少冗余。更新效果取决于模型对信息增量和轨迹质量的判断。

推理时，系统依据任务和截图寻找相似种子，再沿图扩展邻居并按相似度排序。主实验每次使用10条轨迹，包括5条语义相似种子和从其1跳邻居中选出的5条轨迹。一般算法允许进一步扩展邻居。

\subsection{混合工作记忆及局部刷新}

检索后形成两个互补视图：策略与属性归纳为指导指令（Guidance Instructions），作为文本写入系统提示；轨迹连续嵌入则与模型输入连接。前者帮助规划，后者提供细粒度视觉和动作信息，二者构成当前任务的工作记忆。

每个动作后，视觉语言模型比较前后观测与当前指导，判断是否发生阶段转换。若任务由搜索转入结账，先保留长期目标和约束，丢弃不相关内容，再按新状态检索，合并指导并刷新连续嵌入。这是\textbf{任务内部的上下文更新}，与长期图的ADD/MERGE/REPLACE属于不同层次。

\subsection{实验设置与结果}

实验覆盖WebVoyager、Multimodal-Mind2Web和MMInA三个多模态网页智能体基准。前两者根据任务描述与截图轨迹进行模型判分；MMInA比较最终回答和标准答案。主指标是任务成功率，Overall按可用领域取平均。各方法在统一环境、抽样和评价方式下实现。

记忆来源为GUIAct、Mind2Web训练集及作者智能体生成的成功轨迹，共2883条；主实验图含1858个节点和超过100万条边。执行采用ReAct式接口，包含SOM界面标记及UI-INS-7B定位回退。训练更新记忆编码器的Q-Former和LoRA层，可训练参数约占总参数的1.2\%，通常训练1个epoch。系统效果因此同时依赖记忆数据、编码器训练和界面定位能力。

\begin{table}[!htbp]
\centering
\caption{不同记忆表示下的Overall任务成功率（\%）}
\label{tab:hymem-results}
\begin{tabular}{lrrrr}
\toprule
基础模型 & 无记忆 & 离散记忆 & 连续记忆 & HyMEM \\
\midrule
Qwen2.5-VL-7B & 12.5 & 17.0 & 21.7 & 35.0 \\
Qwen3-VL-8B & 21.9 & 21.3 & 23.3 & 33.8 \\
UI-TARS-1.5-7B & 17.1 & 15.4 & 28.7 & 31.2 \\
\bottomrule
\end{tabular}
\end{table}

Qwen2.5-VL-7B由12.5\%升至35.0\%，提高22.5个百分点，混合记忆高于单独的离散或连续记忆。不同模型和领域的收益有差异：Qwen3-VL-8B离散记忆的21.3\%略低于无记忆的21.9\%，部分领域采用HyMEM后也未改善。

在Amazon与Google Maps的更新分析中，全局演化使成功率分别由48.8\%升至61.0\%、43.9\%升至51.2\%；局部刷新分别达到56.1\%和48.8\%。两项分析各自使用静态记忆作为参照，展示跨任务更新和任务内刷新的作用。

\subsection{方法局限与课设启发}

HyMEM结合可解释策略与连续视觉细节，并通过长期图更新和任务内刷新适应状态变化。其开销包括记忆构建、编码器训练、图扩展和执行期间的模型判断。图关联与更新依赖标签质量及模型判断，错误记忆传播、阶段误判和跨网站迁移是进一步评测的重点。

对文本事件记忆而言，可借鉴HyMEM对语义种子与结构邻居、长期更新与工作记忆刷新的区分。课程设计可在固定检索预算下比较邻居扩展带来的证据覆盖、无关信息和调用成本；连续轨迹编码则适用于具有截图和动作序列的多模态任务。

% 原文定位：PDF第1--3页，表1；PDF第3--5页，\S3.1--3.3、图1；PDF第3页，式(1)；PDF第4页，图1与Memory Update；PDF第4--5页，Structured Retrieval及Memory Settings；PDF第5页，On-the-fly Working Memory Refresh；PDF第5--7页，\S4.1--4.2与表2；PDF第5页，\S4.1；PDF第7页，\S4.3.1、图2。
